\section{Introduction}

We study the exceptional set estimate and a Furstenberg-type problem in prime fields. The main novelty of this paper is that we prove, for these two problems, that the two-dimensional result implies all higher-dimensional results. 

\subsection{Background}

We first introduce the two problems---the Furstenberg set problem and the exceptional set estimate in Euclidean space.




\subsubsection{Furstenberg set problem}\hfill

The Furstenberg set problem was originally considered by Wolff \cite{wolff1999recent} in $\R^2$. Fix $s\in (0,1], t\in (0,2]$. We say $E\subset \R^2$ is an $(s,t)$-set in $\R^2$, if $E$ is of the form
\[ E=\bigcup_{\ell\in \cL}Y(\ell), \]
where $\cL\subset A(1,\R^2)$ is a set of lines in $\R^2$ with $\dim\cL=t$ and for each $\ell\in \cL$ there exist $Y(\ell)\subset \ell$ with $\dim Y(\ell)=s$. ($\dim$ will always mean Hausdorff dimension.) The Furstenberg set problem asks about the optimal lower bound on the Hausdorff dimension of $(s,t)$-sets, that is,
\[ \inf_{E \textup{~is~a~}(s,t)\textup{-set}}\dim(E). \]
Recently, Ren and Wang \cite{ren2023furstenberg} resolved the Furstenberg set problem by proving
\begin{equation}\label{renwang}
\inf_{E \textup{~is~a~}(s,t)\textup{-set}}\dim(E)=\min\{s+t,\frac{3}{2}s+\frac{1}{2}t,s+1  \}.
\end{equation} 
For more references and history on this problem, we refer to \cite{ren2023furstenberg}, \cite{orponen2023projections}. 
